\section{Implementazione}
\label{sec:implementazione}

\subsection{Aggregazione Edge e Cloud}

L'elaborazione utilizza finestre tumbling basate sull'event time. Ciascun
Edge mantiene una finestra corrente e assegna ogni evento alla finestra
determinata dal relativo timestamp. Per ogni grandezza ambientale vengono
mantenuti il numero di valori validi e non validi, la somma, il minimo e il
massimo. Al momento della chiusura della finestra viene inoltre calcolata la
media sui soli valori validi.

Le statistiche sono state scelte in modo da poter essere combinate nei
livelli successivi senza conservare le singole osservazioni. Indicando con
$n_j$ il numero di valori validi e con $s_j$ la loro somma nel contributo
$j$, la combinazione di più aggregati utilizza

\begin{equation}
    n = \sum_j n_j,
    \qquad
    s = \sum_j s_j,
    \qquad
    \bar{x} = \frac{s}{n}.
\end{equation}

Minimo e massimo globali sono ottenuti rispettivamente come minimo e massimo
dei contributi che contengono almeno un valore valido. In questo modo la media
globale non viene calcolata come media delle medie locali, che produrrebbe un
risultato errato quando gli aggregati contengono numerosità differenti.

Le finestre Cloud hanno durata configurabile, ma devono essere un multiplo
intero della durata delle finestre Edge. Ogni \code{EdgeAggregate} appartiene
quindi interamente a una sola finestra Cloud. I Cloud Worker combinano gli
aggregati ricevuti per partizione Kafka e producono un
\code{CloudPartitionAggregate}; il Global Aggregator applica infine la stessa
operazione di riduzione ai contributi delle diverse partizioni relativi allo
stesso intervallo temporale.

\subsection{Progresso event-time e completamento}

La chiusura delle finestre Cloud è governata dal progresso in event time
delle sorgenti e non da timeout in processing time. Ogni
\code{EdgeAggregate} contiene il campo \code{CompleteThrough}, che certifica
che l'Edge non produrrà successivamente aggregati con fine della finestra
precedente o uguale a tale istante.

Quando l'arrivo di un evento determina il passaggio a una nuova finestra,
l'Edge emette l'aggregato della finestra precedente e può avanzare
\code{CompleteThrough} fino all'inizio della nuova finestra. In questo modo
vengono certificate anche eventuali finestre intermedie prive di eventi.
Gli Edge non producono infatti aggregati espliciti per le finestre vuote.

Il Cloud mantiene separatamente il progresso di ciascun Edge appartenente
alla partizione Kafka elaborata. Indicando con $C_e$ il valore
\code{CompleteThrough} dell'Edge attivo $e$ e con $S$ il massimo disallineamento
ammesso tra le sorgenti, la frontiera della partizione è determinata da

\begin{equation}
W_P =
\max\left(
    \min_{e \in E_P} C_e,\;
    \max_{e \in E_P} C_e - S
\right).
\end{equation}

Finché almeno un Edge attivo non ha ancora comunicato il proprio progresso,
la frontiera non può avanzare. In condizioni normali essa coincide quindi
con il minimo progresso degli Edge della partizione; il secondo termine
impedisce invece che una sorgente eccessivamente in ritardo blocchi
indefinitamente la finalizzazione delle finestre. Una finestra Cloud viene
emessa quando il proprio estremo superiore non è successivo a $W_P$.

Un aggregato relativo a una finestra già superata dalla frontiera è
considerato \emph{late}: non modifica né riapre il risultato già emesso.
Il progresso contenuto nel record può comunque essere utilizzato per
aggiornare lo stato dell'Edge e consentirgli di riallinearsi alle altre
sorgenti.

Poiché il workload è ottenuto dal replay di un dataset finito, è inoltre
necessario determinare esplicitamente la conclusione del flusso. Al termine
del proprio replay ciascun Simulator invia un marker di fine input al
rispettivo Edge. L'Edge pubblica l'ultima finestra ancora aperta e,
successivamente, un \code{EdgeEndOfInput}. Quando tutti gli Edge appartenenti
a una partizione sono terminati, il Cloud Worker emette gli eventuali
aggregati ancora pendenti e infine un \code{PartitionEndOfReplay}.

Il Global Aggregator utilizza i risultati delle partizioni insieme ai relativi
progressi e marker terminali per distinguere un contributo non ancora arrivato
da un contributo effettivamente assente. Una finestra globale può quindi
essere prodotta quando, per ogni partizione sorgente, è disponibile il
relativo aggregato oppure è stato certificato che non ne verranno prodotti
altri per quell'intervallo. Il replay termina quando tutte le partizioni hanno
raggiunto lo stato terminale e non rimangono finestre pendenti.

\subsection{Comunicazione e affidabilità}

La comunicazione tra Simulator ed Edge utilizza MQTT. Gli eventi dei sensori
sono serializzati in JSON e pubblicati con QoS~0. JSON è stato scelto in
questo tratto della pipeline per la semplicità della rappresentazione e per
la facilità con cui i messaggi possono essere ispezionati durante lo sviluppo
e il debugging. La telemetria viene inoltre prodotta con frequenza elevata:
l'utilizzo di QoS~0 evita di attendere una conferma per ogni singola
osservazione e permette al replay di procedere secondo il rate configurato.
Il marker di fine replay utilizza invece QoS~1 e viene considerato completato
soltanto dopo la ricezione del relativo PUBACK, poiché la sua perdita
impedirebbe di determinare correttamente la terminazione della sorgente.

La comunicazione tra Edge e Cloud e tra Cloud Worker e Global Aggregator è
invece affidata ad Apache Kafka. In questo tratto i messaggi sono serializzati
in formato Apache Avro. Rispetto al JSON, Avro fornisce una rappresentazione
binaria più compatta e definisce esplicitamente la struttura e i tipi dei
record scambiati, caratteristica utile per aggregati composti da numerosi
campi numerici e temporali. Gli schemi Avro sono inclusi direttamente
nell'applicazione e devono quindi essere condivisi in modo coerente da
producer e consumer; non viene utilizzato uno Schema Registry.

Il tipo logico di ciascun record Kafka è indicato mediante l'header
\code{record\_type}. In questo modo lo stesso meccanismo di trasporto può
veicolare aggregati, informazioni di progresso e marker terminali. Questi
ultimi non richiedono dati applicativi e vengono pertanto inviati con payload
vuoto.

Gli Edge utilizzano writer Kafka sincroni con acknowledgment del broker. Gli
aggregati appartenenti allo stesso Edge e il relativo
\code{EdgeEndOfInput} sono scritti sulla medesima partizione; prima di
pubblicare il marker terminale vengono completate le scritture degli aggregati
precedenti. Analogamente, il Cloud pubblica per ciascuna partizione prima gli
eventuali \code{CloudPartitionAggregate}, quindi il relativo
\code{PartitionProgress} e infine, quando necessario,
\code{PartitionEndOfReplay}. L'ordinamento permette ai consumer downstream di
interpretare un'informazione di progresso solo dopo aver ricevuto tutti i dati
che essa certifica.

I Cloud Worker effettuano il commit degli offset Kafka solamente dopo che il
record di input è stato elaborato e gli eventuali output sono stati pubblicati
con successo. I commit possono essere raggruppati in batch configurabili per
ridurne il costo; la ricezione di un marker terminale forza invece il commit
degli offset ancora pendenti. Il Global Aggregator applica lo stesso principio,
effettuando il commit dopo l'elaborazione e la scrittura del risultato nel
sink.

Questa strategia evita di confermare un record prima che i suoi effetti siano
stati propagati, ma non realizza una transazione atomica end-to-end tra stato,
output Kafka e offset di consumo. Un arresto dopo la pubblicazione dell'output
ma prima del commit può quindi causare la rilettura di un record. Per questo
motivo gli aggregati utilizzano identificativi deterministici e i componenti
riconoscono i duplicati identici, mentre nel sink PostgreSQL l'identificativo
dell'aggregato è una chiave univoca.

\subsection{Deployment e piattaforma software}

L'applicazione è sviluppata in Go e ciascun componente principale viene
distribuito come container Docker. La configurazione del sistema è descritta
tramite file YAML, dai quali il programma \code{deploygen} genera
automaticamente i file Docker Compose necessari all'esecuzione. In questo modo
parametri quali numero di Cloud Worker, numero di partizioni Kafka, dimensione
delle finestre e parametri di batching possono essere modificati senza
intervenire sul codice applicativo.

Per il deployment distribuito l'infrastruttura AWS è definita mediante
Terraform. Sono previsti quattro host EC2 con ruoli distinti: Simulator, Edge,
Cloud Core e Workers. L'host Simulator esegue i processi di replay, l'host Edge
ospita i broker Mosquitto e gli Edge Gateway, l'host Workers esegue le repliche
dei Cloud Worker, mentre Cloud Core ospita Kafka e il Global Aggregator.
La separazione consente di far attraversare realmente la rete ai flussi
Simulator--Edge ed Edge--Cloud e di modificare indipendentemente il numero di
worker utilizzato negli esperimenti.

La persistenza dei risultati globali può essere affidata a un'istanza
PostgreSQL gestita tramite Amazon RDS. Il database non è esposto pubblicamente
ed è raggiungibile dal solo Cloud Core. L'infrastruttura, la configurazione
applicativa e l'avvio degli esperimenti sono gestiti tramite script dedicati,
in modo da rendere ripetibile la preparazione dell'ambiente e la raccolta degli
artefatti prodotti da ogni esecuzione.

Le principali librerie utilizzate dall'applicazione sono Eclipse Paho per MQTT,
\code{kafka-go} per la comunicazione con Kafka, \code{goavro} per la
serializzazione Avro e \code{pgx} per l'accesso a PostgreSQL. Il preprocessing
offline del dataset è invece realizzato in Python mediante pandas,
scikit-learn e pyproj. Non viene utilizzato un framework esterno di stream
processing: la gestione delle finestre, dello stato e del progresso event-time
è implementata direttamente nei componenti Go.